\clearpage
\chapter*{Colophon: archiving these notes}
\addcontentsline{toc}{chapter}{Colophon: archiving these notes}
\label{ch:colophon}

These notes are a worked example of their own subject. Their \LaTeX{} source is archived in
Software Heritage, and the intrinsic identifier of the archived source directory is:

\begin{center}
\ifdefempty{\thisswhid}{%
  \emph{(stamped at release time --- see ``How this identifier is produced'' below)}%
}{%
  \swhid{\thisswhid}%
}
\end{center}

\paragraph{Why this identifier is not written into the source.}
A \textsc{swhid} is a cryptographic hash of the content it names. Writing a document's own
\textsc{swhid} \emph{into} that document would change the document's bytes, and therefore
change its hash --- and you cannot, in any case, embed a digest you have not computed yet, and
cannot compute until the bytes are fixed. This is not a limitation of Software Heritage; it is what
makes the identifier
trustworthy: the name is forced to follow the bytes, so it cannot be faked. The source is
therefore archived first, and the resulting identifier is stamped only into \emph{this
rendered PDF}, which is not part of the hashed source tree.

\paragraph{How this identifier is produced.}
After the source is committed and archived (via Save Code Now), the qualified directory
\textsc{swhid} is read from the archive's \emph{Permalinks} tab and stamped into a
non-committed PDF:
\begin{center}
\texttt{make swhid SWHID='swh:1:dir:...;origin=...;visit=...'}
\end{center}
A plain \texttt{make} deliberately leaves the slot empty: the committed source never contains
its own identifier, so the source's \textsc{swhid} stays well-defined and stable.

\paragraph{If you obtained this PDF from HAL (or any cover-page repository).}
Repositories such as HAL prepend an institutional cover page to every deposited file, and may
re-render or re-compress it. That copy is therefore a \emph{different digital object}: its bytes
--- and any content hash taken over them, such as a \swhid{swh:1:cnt:} --- will not match this
rendering. The identifier above is unaffected, because it names the archived \emph{source
directory}, not the PDF: a source tree has a canonical form, so the capstone below verifies from
any copy. But the pristine, canonical rendering of this note lives at
\begin{center}\url{https://dicosmo.org/Articles/2026-source-code-of-science.pdf}\end{center}
and that is the copy to hash, diff, or byte-archive if you need bytes that match.

\paragraph{Why source code gets a perfect identifier and a PDF does not.}
This caveat is the whole note in miniature. An intrinsic identifier pins \emph{source code}
exactly \emph{because source has a canonical form}: the same tree always hashes to the same
\textsc{swhid}. A rendered document has no single canonical byte-form --- a prepended cover page,
a newer \TeX{} Live, re-embedded fonts, or a re-compression each produce a perfectly valid PDF
with a \emph{different} hash. So identify the \emph{source} by its \texttt{dir} \textsc{swhid},
which is stable and reproducible, and treat any PDF --- this one included --- as one rendering of
it rather than as the thing itself.

\begin{handson}[Hands-on (capstone): verify this very note]
\labmeta{$\sim$10 min \textperiodcentered{} needs git $+$ a \textsc{swhid} tool \textperiodcentered{} in-class, offline}
You can recompute the identifier above and confirm it yourself.
\begin{enumerate}[nosep,leftmargin=*]
  \item Export the \emph{committed} source (tracked files only, no build artifacts):\\
        \texttt{git archive HEAD | (mkdir -p /tmp/note \&\& tar -x -C /tmp/note)}
  \item Compute the directory \textsc{swhid} with the Rust reference tool (\texttt{cargo install
        swhid}; see Appendix~\ref{app:swhid-tools}): \texttt{swhid dir /tmp/note}.
  \item Compare its core \swhid{swh:1:dir:<hash>} with the one printed above (ignore the
        \texttt{origin}/\texttt{visit} qualifiers). A match means you have verified the
        integrity of the document you are reading.
  \item \emph{Bonus (git-compatibility):} the notes' \emph{revision} identifier is just
        \swhid{swh:1:rev:} prepended to \texttt{git rev-parse HEAD} --- a Software Heritage
        revision \textsc{swhid} \emph{is} the git commit hash.
\end{enumerate}
\end{handson}
